\section{Zero geometry, exhaustion, and counting}
\subsection{Local zeros are only the start}
The two-term dominance analysis separates the leading pair of exponentials from a controlled remainder. Near the cancellation lattice of the leading pair, the perturbation argument constructs nearby zeros and estimates their displacement. The remainder bounds include derivative information, which supports simplicity and the first correction term.

A local construction alone would not justify counting all zeros. The source therefore adds sector confinement, sector exhaustion, rotation, and global exhaustion. Those modules show that sufficiently large zeros belong to the constructed symmetry strings. Finite exceptional zeros remain present throughout the later counting and product arguments.

The following division of labor is useful when reviewing the proof path:
\begin{longtable}{p{.19\textwidth}p{.48\textwidth}p{.23\textwidth}}
\toprule Stage & Output & Map rows \\ \midrule\endhead
Remainder control & Canonical polygon gap, strip and disk bounds, normalized jets & 58--60, 67 \\
Local perturbation & Nearby zeros, displacement, derivative estimates & 57, 61 \\
Global coverage & Sector confinement, exhaustion, rotational orbits & 62--65 \\
Simplicity and rays & All sufficiently large zeros simple and on exact rays & 66, 68 \\
Reduced variable & Zero and multiplicity transfer under $z\mapsto z^N$ & 69--71 \\
Complete enumeration & Constructed tail plus finite exceptions & 75--79, 81 \\
Refined phase & Real leading correction and alternating derivative sign & 82--84 \\
\bottomrule
\end{longtable}

\subsection{Named enumeration is constructed, not assumed}
The \code{SignedZeroEnumeration} structure packages a positive simple tail, its lattice windows, rotation orbits, and the finite-complement information needed downstream. Existence of this structure is proved for every $N\geq2$, with the exact order-two cosine model handled separately.

A downstream theorem may accept an enumeration as input because it is reusable for any structure with the required fields. That is different from assuming the final zero theorem. The manuscript specializations instantiate it using \code{canonicalZeroEnumeration}, constructed from the preceding analytic results. To audit a final theorem, follow that constructor as well as the generic lemma.

\subsection{Multiplicity enters the actual counting function}
The disk counting function sums \code{analyticOrderNatAt} over the actual finite zero set. The source proves that these analytic orders are finite for the nonzero entire denominators. Thus replacing an extended order by its natural-number value does not silently discard an infinite order.

The full original-variable counting theorem keeps every exceptional multiplicity:
\excerpt{FullZeroCounting}{counting}
The proof compares the full count with the simple tail count and adds a finite exception weight. The final error constant contains both that weight and the tail shift. It is uniform for all nonnegative radii, not only an eventual asymptotic abbreviation.

The corresponding reduced-variable count is
\[
 n_{\Delta_N}(R)=\frac{\sin(\pi/N)}{\pi}R^{1/N}+O(1),\qquad N\geq2.
\]
The original-variable constant is $N\sin(\pi/N)/\pi$. These constants differ because one original symmetry orbit collapses to one reduced zero and the radius changes by the $N$th power. Both transformations are proved in the counting modules.

\subsection{Sharp summability thresholds}
Rows 72 and 79 distinguish a tail estimate from the full theorem. The full reciprocal-power sums range over actual zeros with their analytic multiplicities. Their sharp thresholds are $s>1$ for the original denominator and $s>1/N$ for the reduced denominator. The boundary divergences are included. Removing finite exceptions is justified by a proved summability equivalence, not by silently redefining the zero set.

For $N\geq2$, reduced reciprocal zeros are absolutely summable because $1>1/N$. This precise consequence feeds the genus-zero product and the trace-class diagonal model. It would fail as a route at $N=1$, where the denominator is exponential and the analytic situation is handled separately.

\subsection{What eventual simplicity leaves open}
The source proves eventual simplicity and exact rays for sufficiently large zeros. It does not replace this with a blanket assertion that all zeros are simple. A finite exceptional zero can have multiplicity greater than one, and the canonical product includes it accordingly. A finite exceptional zero can also lie outside the asymptotic positive-ray description. The global product and zero-zeta constructions are designed to retain such exceptions.

The first displacement correction and its exponential error are likewise specific proved statements. They do not constitute an all-orders expansion. The explicit cubic correction in row 84 is a specialization with exact spectral gap, lattice centers, and phase, not a numerical fit to plotted zeros.
